\section{Conclusion}\label{sec:conclusion}

We have demonstrated the utility of representing policies using text-conditioned video generation, showing that this enables effective combinatorial generalization, multi-task learning, and real world transfer. These positive results point to the broader direction of using generative models and the wealth of data on the internet as powerful tools to generate general-purpose decision making systems.

\myparagraph{Limitations.} Our current approach has several limitations. First, the underlying video diffusion process can be slow, it can take a minute to generate highly photorealistic videos. This slowness can be overcome by distilling the diffusion process into a faster sampling network~\citep{salimans2022progressive}, which in our initial experimentation resulted in a 16x speed-up. UniPi may further be sped up with by faster speed samplers in diffusion models. Second, the environments considered in this work are generally fully observed. In partially observable environments, video diffusion models might make hallucination of objects or movements that are unfaithful or not in the physical world. Integrating video models with semantic knowledge about the world may help resolve this issue, and the integration of UniPi with LLMs would be an interesting direction of future work.

